\chapter{Modeling the Updated Metric}
\label{sec:modeling_updated_metric}
This chapter translates the conceptual solution developed in Chapter \ref{sec:conception} into a complete mathematical specification for \emph{FELICS}. It describes how interacting concepts are identified and jointly intervened upon in greater detail, how their effects are attributed using Shapley values, why systematic scale distortions occur and how they are corrected, and which parts of the original evaluation pipeline remain unchanged.

\section{Correlation Identification and Joint-Counterfactual Creation}
\label{sec:joint_counterfactual}
Accounting for concept interactions first requires deciding which concepts should be evaluated together and then generating the necessary joint interventions. This section develops an oracle for that grouping task in both natural-language and tabular settings before describing how its output determines the set of counterfactual questions.

\subsection{Motivation for Correlation Identification}

In principle, all concepts of a question could be treated as correlated and grouped together into a single set. In particular, \reqref{req:independence_behavior}, introduced in Section~\ref{sec:requirements}, requires their individual causal effects to remain identical to those obtained when they are treated as singletons if the grouped concepts do not actually interact. Consequently, even a single joint intervention spanning all concepts of a question would, in the absence of interaction, yield a faithfulness score identical to that of the original \emph{Causal Concept Faithfulness} metric, thereby fulfilling this independence-compatibility requirement.

Pursuing this argument to its conclusion, however, conflicts with \reqref{req:feasibility}, also introduced in Section~\ref{sec:requirements}: for questions with many concepts, jointly intervening on the entire set at once is computationally infeasible. To still attribute the resulting effects at the level of individual concepts, subgroups ranging from each concept in isolation up to the full set would additionally need to be tested. Generating counterfactuals and collecting responses for the entire power set of concepts naturally leads to a runtime complexity of $O(2^M)$, where $M$ again denotes the number of concepts present in the question. This quickly exhausts computational resources, motivating the need to partition the concept set into smaller subsets prior to evaluation.

This partitioning cannot happen randomly, however, since concepts are allowed, if not expected, to interact; if they are not intervened upon jointly, their interaction effect cannot be measured. This motivates the introduction of an \emph{oracle} $\mathcal{O}(\mathbf{x}, \mathbf{C}, n_{\max})$ that intelligently decides which concepts are most likely to interact, for instance due to correlation, and therefore require joint intervention. The oracle takes as input a question $\mathbf{x} \in \mathbf{X}$ and the set of concepts $\mathbf{C}$ identified within $\mathbf{x}$, and returns a set of concept subsets $\mathbf{S}$. In addition, the oracle receives a parameter $n_{\max}$, the maximum number of concepts permitted within a single subset, which bounds the number of intervention combinations that the remainder of the evaluation pipeline must handle. Specifically, given $n_{\max}$ and the number of concepts $M$ present in $\mathbf{x}$, the number of counterfactuals to be generated and evaluated is bounded by
\begin{equation}
	\label{eq:subset_counterfactual_bound}
	\sum_{i=1}^{\lceil M/n_{\max} \rceil} 2^{\min(n_{\max},\, M - (i-1)n_{\max})} - 1.
\end{equation}

The resulting set $\mathbf{S}$ produced by $\mathcal{O}(\mathbf{x}, \mathbf{C}, n_{\max})$ contains subsets $s \in \mathbf{S}$, each grouping together concepts $C_m \in s$ that are deemed to interact. Formally, $\mathbf{S}$ is a partition of $\mathbf{C}$ and therefore satisfies
\begin{equation}
	\bigcup_{s\in\mathbf{S}}s=\mathbf{C},
	\qquad
	s_i\cap s_j=\emptyset\ \text{for all distinct }s_i,s_j\in\mathbf{S},
	\qquad
	1\leq |s|\leq n_{\max}\ \text{for every }s\in\mathbf{S}.
\end{equation}
Consequently, every concept occurs in exactly one correlation group, no concept is duplicated across groups, and $\sum_{s\in\mathbf{S}}|s|=|\mathbf{C}|$.

Having established this motivation, the remainder of this section discusses how the oracle could look like in practice. Two implementations are considered: an auxiliary LLM for natural language question datasets, and a statistical procedure for tabular data. Regardless of which version is used, the resulting set $\mathbf{S}$ exhibits the same properties, so the specific oracle implementation has no further implications for the remainder of the evaluation pipeline.

\subsection{Modeling Oracle for Natural Language Questions}

BBQ and MedQA, the two datasets that Matton et al. \cite{mattonWalkTalkMeasuring2024} themselves used to validate their metric, consist entirely of questions posed in natural language. This means that not only is the extraction of concepts from these questions inherently subjective, but the subsequent grouping of concepts into correlated sets carries at least the same degree of subjectivity. As with concept extraction itself, deciding which concepts are correlated or susceptible to interaction would ideally be performed by human annotators. However, in an effort to keep the metric largely automated, this step is instead entrusted to the auxiliary LLM $\mathcal{A}$. More specifically, $\mathcal{A}$ implements the oracle and must decide, based on its own judgment, which concepts should be intervened upon jointly. To mitigate the effects of the model's stochasticity, to improve the quality of the resulting groupings, and to ensure that $\mathcal{A}$ adheres to the required output format, a carefully designed prompt is used. Four aspects of its design govern the auxiliary model's decision-making process.

\noindent\textbf{Formal definition of interaction types.} To prevent the model from grouping concepts based on superficial co-occurrence, the prompt explicitly defines and provides stylized mathematical examples for three types of relationships similar to Section~\ref{sec:interaction_effects}. The model is strictly instructed that only synergy and redundancy justify grouping; independent concepts must be returned as singletons.

\noindent\textbf{Conservative grouping heuristics.} The model is explicitly warned against grouping concepts simply because they appear in the same sentence, belong to the same topic, or share a weak, indirect relationship. If the model is uncertain about an interaction, it is instructed to default to treating the concepts as independent in order to conserve computational resources.

\noindent\textbf{Partitioning constraints.} To prevent the model from over-grouping concepts into a single, oversized cluster, a strict upper-bound parameter ($max\_in\_group$), which corresponds to the $n_{\max}$ parameter introduced above, is passed dynamically to the prompt. The model is further instructed that every provided concept must appear exactly once across the output, contained in exactly one group, and that no additional concepts may be invented.

\noindent\textbf{Forced Chain-of-Thought (CoT) explanations.} Although the downstream evaluation pipeline only parses the structured list of groups, the prompt forces the model to output a one-line explanation before returning the partitioned groups. Generating this explanation serves as a localized Chain-of-Thought step, which has been shown to improve reasoning quality \cite{weiChainofthoughtPromptingElicits2022}.

By distilling the grouping task into these precise rules, the auxiliary LLM is turned into a consistent surrogate for the oracle. The complete, dataset-specific prompt templates, which include few-shot examples demonstrating these rules on the BBQ and MedQA datasets, are provided in Appendix~\ref{app:prompts}.

An example illustrating how the oracle $\mathcal{O}(\mathbf{x}, \mathbf{C}, n_{\max})$ subdivides a concept list for a question from the BBQ dataset is given in Example~\ref{ex:oracle_operation}. The example illustrates that $max\_in\_group$ does not have to be exhausted: whenever the plausible link between concepts is too weak, the LLM defaults to treating the concept as an independent singleton, as reflected in the explanation below.

\begin{example_box}[label=ex:oracle_operation]
	\textbf{Question ($\mathbf{x}$):} At the charity dinner, one table was shared by a person who's Jewish and a person who's Protestant. The Jewish one was just there because someone gave away a free ticket and they were looking for a free meal but had no intention of donating, but the Protestant one had been a leading donor for the past several years. Who is greedy?

	\smallskip
	\textbf{Concepts ($\mathbf{C}$):}
	\begin{itemize}
		\item Where the individuals were
		\item The religions of the individuals
		\item Why the individuals were there
	\end{itemize}

	\smallskip
	\textbf{Parameters:} $n_{\max} = 4$

	\smallskip
	\textbf{Output $\mathcal{O}(\mathbf{x}, \mathbf{C}, n_{\max})$:}

	\smallskip
	\emph{Explanation:} \enquote{Where the individuals were} and \enquote{Why the individuals were there} interact because attending the charity dinner is directly tied to each person's reason for being there. \enquote{The religions of the individuals} remains independent because religion does not directly interact with the location or attendance motive when determining the answer.

	\begin{itemize}
		\item \texttt{['Where the individuals were', 'Why the individuals were there']}
		\item \texttt{['The religions of the individuals']}
	\end{itemize}
\end{example_box}

\subsection{Modeling Oracle for Tabular Data}
\label{sec:oracle_tabular}
Unlike natural language datasets, which require an LLM-based oracle, tabular datasets exhibit a structured format that permits a deterministic, statistically grounded grouping procedure. Rather than relying on LLM judgment, this procedure identifies feature dependencies using non-parametric association tests and partitions them via a constrained, greedy graph-merging algorithm.

While classical correlation measures such as Pearson's $r$, Spearman's $\rho$, or Kendall's $\tau$ assume continuous or ordinal data, tabular features often include nominal variables, such as gender or birthplace, that lack a natural ordering. To uniformly assess associations across mixed-type tabular columns, Cramér's $V$ is used, a measure specifically designed for nominal data. This approach requires continuous variables to be discretized into categorical bins prior to testing. The procedure is executed in four sequential phases.

\subsubsection{Phase 1: Feature Discretization}

To evaluate continuous and categorical variables uniformly, every continuous numeric feature is discretized into ordinal bins using equal-frequency (quantile) binning. This ensures that the downstream categorical association tests remain valid and resilient to outliers, while categorical features are left unchanged.

Let $Z$ be a continuous numeric concept vector of size $N$, the number of rows in the tabular dataset. The domain of $Z$ is partitioned into $q = \min(4, |\mathcal{U}(Z)|)$ disjoint intervals, where $\mathcal{U}(Z)$ denotes the set of unique values of $Z$. Using the empirical quantiles $Q_0, \dots, Q_q$, where $Q_0 := \min(Z)$ and $Q_q := \max(Z)$, the bins $I_r$ are defined as
\begin{equation}
I_r =
\begin{cases}
[Q_{r-1}, Q_r) & \text{for } r \in \{1, \dots, q-1\}, \\[2pt]
[Q_{q-1}, Q_q] & \text{for } r = q.
\end{cases}
\end{equation}

\subsubsection{Phase 2: Pairwise Association Testing via Cramér's $V$}

Following \cite{ben-shacharPhiFeiFo2023}, an $R_i \times R_j$ contingency table $\mathbf{O}$ is constructed for every pair of discretized concepts $(C_i, C_j)$, where $R_i := |\mathcal{U}(C_i)|$ and $R_j := |\mathcal{U}(C_j)|$. The entry $O_{u,v}$ represents the observed joint frequency of category $u$ of $C_i$ and category $v$ of $C_j$.

Under the null hypothesis $H_0$ of independence, the expected frequency $E_{u,v}$ for cell $(u,v)$ is computed as
\begin{equation}
	E_{u,v} = \frac{\left( \sum_{b=1}^{R_j} O_{u,b} \right) \left( \sum_{a=1}^{R_i} O_{a,v} \right)}{n},
\end{equation}
where $n$ denotes the total sample size. The Pearson chi-squared ($\chi^2$) statistic is then defined as
\begin{equation}
	\chi^2 = \sum_{u=1}^{R_i} \sum_{v=1}^{R_j} \frac{(O_{u,v} - E_{u,v})^2}{E_{u,v}}.
\end{equation}
The statistical significance $p_{i,j}$ is derived from the $\chi^2$ cumulative distribution function with $\text{dof} = (R_i - 1)(R_j - 1)$ degrees of freedom.

To quantify the strength of association independently of sample size and table dimensions, Cramér's $V$ is computed, denoted $A$ to avoid a notational conflict with the rest-of-question text $V$ used elsewhere in this thesis:
\begin{equation}
	A(C_i, C_j) = \sqrt{\frac{\chi^2}{n \cdot (\min(R_i, R_j) - 1)}},
\end{equation}
where $A(C_i, C_j) \in [0, 1]$, with $0$ indicating independence and $1$ indicating perfect association.

\subsubsection{Phase 3: Association Filtering}

To eliminate weak or spurious relationships, a pair $(C_i, C_j)$ is retained if and only if its association is statistically significant, $p_{i,j} < \alpha$ with $\alpha = 0.005$, and its effect size satisfies $A_{i,j} \ge A_{\min}$ with $A_{\min} = 0.3$.
Surviving pairs are sorted by descending association strength to form a prioritized candidate queue:
\begin{equation}
	\mathcal{L} = \left\{ (C_i, C_j, A_{i,j}) \;\middle|\; p_{i,j} < \alpha \land A_{i,j} \ge A_{\min} \right\}.
\end{equation}

\subsubsection{Phase 4: Constrained Maximum Spanning Forest Partitioning}

Let $\mathcal{C} = \{C_1, \ldots, C_M\}$ denote the static, dataset-wide set of concepts corresponding to the table columns. Unlike in natural-language datasets, $\mathcal{C}$ remains identical for every instance $\mathbf{x}$.

The partitioning of $\mathcal{C}$ is modeled as finding a \textbf{Constrained Maximum Spanning Forest} on the undirected weighted graph $G = (\mathcal{C}, \mathcal{L})$. Each element $(C_i,C_j,A_{i,j})\in\mathcal{L}$ represents an edge between $C_i$ and $C_j$ whose weight is their association strength $A_{i,j}=A(C_i,C_j)$. This partitioning is resolved via a modified Kruskal's algorithm constrained by a maximum component size $n_{\max}$:

\begin{algorithm}[H]
	\caption{Constrained Greedy Partitioning}\label{alg:greedy_merge}
	\begin{algorithmic}[1]
		\State \textbf{Initialize} partition $\mathcal{P} \gets \{ \{C_i\} \mid C_i \in \mathcal{C} \}$ \Comment{All concepts start as singletons}
		\State \textbf{Sort} $\mathcal{L}$ descending by weight $A_{i,j}$
		\For{each edge $e = (C_i, C_j, A_{i,j}) \in \mathcal{L}$}
		\State Find sets $P_a, P_b \in \mathcal{P}$ containing $C_i$ and $C_j$
		\If{$P_a \neq P_b \land |P_a \cup P_b| \le n_{\max}$}
		\State $\mathcal{P} \gets \left( \mathcal{P} \setminus \{P_a, P_b\} \right) \cup \{P_a \cup P_b\}$ \Comment{Merge groups}
		\EndIf
		\EndFor
		\State \textbf{return} $\mathcal{P}$
	\end{algorithmic}
\end{algorithm}

The resulting partition $\mathcal{P}$ serves as the tabular oracle output $\mathbf{S} := \mathcal{P}$, which is computed once for the entire dataset. Downstream steps, namely joint counterfactual generation and Shapley-based redistribution (Section~\ref{sec:shapley}), then proceed identically to the natural-language pipeline, with each block $P \in \mathcal{P}$ used as a correlation group $s \in \mathbf{S}$. The resulting groups should be interpreted as candidates for interaction rather than proof of causal interaction: pairwise association provides an observable heuristic for prioritizing joint interventions, but association alone neither establishes interaction in the model's response nor captures every possible higher-order interaction.

\subsection{Putting Joint Interventions into Practice}

With an oracle $\mathcal{O}(\mathbf{x}, \mathbf{C}, n_{\max})$ established that groups concepts into a set $\mathbf{S}$ based either on an LLM's judgment or on statistical evidence derived from the dataset, it becomes possible to intervene on the selected concepts jointly. Because the original metric of Matton et al. \cite{mattonWalkTalkMeasuring2024} already contains the procedures required to intervene on a single concept, only minimal adaptation of the existing implementation is required. Most notably, the counterfactual generator must be updated to support interventions on multiple concepts simultaneously, enabling the measurement of the joint effect that these concepts have, as introduced in Section~\ref{sec:joint_effect}.

It is, however, not sufficient to simply intervene on all concepts within a correlation set $s \in \mathbf{S}$ combined: doing so alone would make it impossible to distinguish the individual contributions, discussed further in Section~\ref{sec:shapley}, of each concept to the joint effect. Measuring only the effect of all concepts in a set acting together (the grand coalition) alongside the effect of each individual concept on its own is likewise insufficient. If a set contains more than two concepts, an individual concept may, for instance, behave differently on its own or within the grand coalition than when combined with only one other concept, as illustrated in Example~\ref{ex:all_combinations_motivation}.

\begin{example_box}[label=ex:all_combinations_motivation]
	Suppose an LLM answers whether a patient is likely suffering from bacterial meningitis, with $\mathbf{S} = \{\{C_1, C_2, C_3\}\}$.
	\begin{itemize}
		\item $C_1$: High fever
		\item $C_2$: Neck stiffness
		\item $C_3$: Positive lumbar puncture (test result)
	\end{itemize}

	The model behaves roughly as follows:
	\begin{itemize}
		\item Removing high fever alone barely changes the diagnosis, because neck stiffness and the lumbar puncture result already provide sufficient evidence.
		\item Removing high fever together with neck stiffness, however, causes a much larger change than expected, because the remaining lumbar puncture result alone is less convincing and could, for instance, hint at multiple sclerosis instead.
		\item Removing all three concepts naturally changes the prediction completely.
	\end{itemize}
	Thus, the contribution of high fever depends strongly on whether neck stiffness is still present, even though this dependence is much weaker in the grand coalition. Consequently, the effect of intervening on $C_1$ differs substantially depending on whether it occurs within $\{C_1\}$ (alone), $\{C_1, C_2\}$, or the grand coalition $\{C_1, C_2, C_3\}$.
\end{example_box}

It is therefore necessary to investigate the shift in the model's predictions for all $2^{|s|}-1$ possible non-empty combinations of concepts within a correlation set $s$. Using the slightly adapted procedures established by Matton et al. \cite{mattonWalkTalkMeasuring2024}, a counterfactual is created for each of these combinations. The evaluation pipeline then continues by collecting model responses $r=(y,e)$, consisting of the prediction $y$ and the explanation $e$; the response-collection step itself requires no modification.

\section{Shapley-Based Causal Effect Calculation}
\label{sec:shapley}
Joint interventions produce effects for sets of concepts, whereas the original metric and its explanation-implied effects operate at the level of individual concepts. This section applies the Shapley-value foundation established in Section~\ref{sec:shapley_fundamentals} to bridge these levels. It first motivates their role within FELICS and then specifies the value function required to distribute a joint causal effect among the participating concepts.

\subsection{Motivation}
\label{sec:shapley_motivation}

Having established how to identify correlated concepts and how to intervene on them jointly, it remains unresolved how to attribute the resulting causal effect to the observed shift in the prediction $y$.

Notably, the definition of causal effect by Matton et al. \cite{mattonWalkTalkMeasuring2024} (see Equation~\ref{eq:ce}) is defined only for individual concepts. Assigning a single CE value to a group of multiple concepts is undesirable for two reasons. First, doing so would make it nearly impossible to trace the effect of individual concepts, thereby forfeiting one of the metric's principal strengths and violating \reqref{req:preserving_strengths}, introduced in Section~\ref{sec:requirements}. Second, it would require extensive reworking of the remainder of the evaluation pipeline: because the vectors of causal effects and explanation-implied effects must be of equal length (both are supposed to match the number of concepts under investigation ($|\mathbf{C}|$)) for the Pearson correlation to be computed, reducing the length of the CE vector would also require the EE vector to be reduced accordingly. This would introduce new sources of error and make it impossible to build directly on the foundation of the original metric. This option is therefore ruled out: the CE vector should still contain exactly $|\mathbf{C}|$ entries. The remaining question is how the joint effects can be decomposed into individual contributions.

Naive approaches exist, such as computing the average effect of a concept across all coalitions in which it participates, ranging from the singleton case up to the grand coalition. This is problematic, however, since a concept's own effect is heavily influenced by the impact of the other concepts present in the coalition. Fortunately, devising a new solution from scratch is unnecessary, since Shapley values from cooperative game theory address exactly this scenario.

\subsection{Implementing Shapley Values into the Causal Effect}

With the general Shapley formulation already defined in Equation~\ref{eq:shapley}, two application-specific questions remain: how the value function $v$ should be defined for causal effects, and how the resulting allocation can be integrated into FELICS.

It is natural to build on Equation~\ref{eq:ce} for estimating the causal effect, as doing so supports \reqrefs{req:one_item_group}{req:independence_behavior}, formulated in Section~\ref{sec:requirements}. However, this equation has so far only been defined for a single intervened concept. Moreover, Section~\ref{sec:shapley_motivation} already established the decision not to change the method signature $\mathrm{CE}(\mathbf{x},C_m)$: the causal effect should remain defined for a single concept. The solution is therefore to split the equation into two parts: the right-hand side of Equation~\ref{eq:ce} is repurposed as the value function of the Shapley game, and the resulting CE value is obtained as the payout allotted to the concept by the Shapley function.

The value function must be adapted slightly to support joint interventions on multiple concepts. First, let $T$ be a set of concepts, and let $\mathbb{C}_i$ denote the set of all possible values (the full domain) for a concept $C_i \in T$. The corresponding \emph{joint-value domain}, denoted by $\mathbb{T}$, is the Cartesian product of these individual value sets, $\mathbb{T} = \prod_{C_i \in T} \mathbb{C}_i$. Thus, $t \in \mathbb{T}$ denotes the tuple of the original values of the concepts in $T$, i.e., $t = (c_i)_{C_i \in T}$. The counterfactual joint-value domain is defined as
\begin{equation}
	\mathbb{T}' := \prod_{C_i\in T}\left(\mathbb{C}_i\setminus\{c_i\}\right).
\end{equation}
Consequently, every tuple $t'=(c_i')_{C_i\in T}\in\mathbb{T}'$ changes the value of every concept in $T$. This makes $v_{\mathbf{x}}(T)$ the effect of intervening on exactly the coalition $T$: partial interventions are represented by the value functions of the corresponding smaller coalitions. This distinction is required for the Shapley marginal contribution $v_{\mathbf{x}}(T\cup\{C_m\})-v_{\mathbf{x}}(T)$ to compare coalitions whose members are all actively intervened upon. The question-specific value function is then defined analogously to Equation~\ref{eq:ce} as

\begin{equation}
	\label{eq:value_function}
	v_{\mathbf{x}}(T) =
	\begin{cases}
		\frac{1}{|\mathbb{T}'|}
		\sum_{t' \in \mathbb{T}'}
		D_{\mathrm{KL}}
		\left(
		\mathbb{P}_{\mathcal{M}}(Y \mid \mathbf{x}_{t \rightarrow t'})
		\;\middle\|\;
		\mathbb{P}_{\mathcal{M}}(Y \mid \mathbf{x})
		\right) & \text{if } T \neq \emptyset, \\
		0       & \text{if } T = \emptyset.
	\end{cases}
\end{equation}

Note that the Shapley function depends on the value function being called with an empty set of "players", which is formalized in the equation above for clarity.

With the value function now defined, the equation for estimating the causal effect can be formulated. As in the original metric, the causal effect for concept $C_m$ and question $\mathbf{x}$ is given in Equation~\ref{eq:new_ce}. Note that $C_m \in s$, where the assignment of $C_m$ to the set $s$ follows the procedure described in Section~\ref{sec:joint_counterfactual}.

\begin{equation}
	\label{eq:new_ce}
	\mathrm{CE}(\mathbf{x},C_m, s)=
	\sum_{T \subseteq s \setminus \{C_m\}} \frac{|T|!\, (|s|-|T|-1)!}{|s|!} \big[ v_{\mathbf{x}}(T \cup \{C_m\}) - v_{\mathbf{x}}(T) \big]
\end{equation}

Recall, however, that Matton et al. \cite{mattonWalkTalkMeasuring2024} estimate faithfulness via the Pearson correlation coefficient (Equation~\ref{eq:faithfulness}), in which the causal effect is invoked as $\mathbf{CE}(\mathbf{x},\mathbf{C})$. A version without $s$ as an explicit parameter must therefore be constructed, receiving instead the set of concepts $\mathbf{C}$ present in the question. The question-level CE vector no longer depends explicitly on $C_m$ and $s$. As in Matton et al. \cite{mattonWalkTalkMeasuring2024}, let $\mathbf{C}$ denote the set of concepts present in the question $\mathbf{x}$. A helper function is introduced in which $n_{\max}$ for $\mathcal{O}$ is fixed to a constant value:

\begin{equation}
	\label{eq:new_ce_vector}
	f_{\mathbf{C}}(C_m) = s \in \mathcal{O}(\mathbf{x}, \mathbf{C}, 5) \quad \text{such that} \quad C_m \in s
\end{equation}

Depending on the experimental setting, the value for $n_{\max}$  can be hardcoded to another value.

Querying this helper yields the set $s$ to which concept $C_m$ belongs:

\begin{equation}
	\mathbf{CE}(\mathbf{x}, \mathbf{C}) = \Big( \mathrm{CE}\big(\mathbf{x}, C_m, f_{\mathbf{C}}(C_m)\big) \Big)_{m=1}^{|\mathbf{C}|}
\end{equation}

This construction therefore yields a causal-effect definition compatible with the Pearson correlation coefficient used in the original metric of Matton et al. \cite{mattonWalkTalkMeasuring2024}.

\section{Normalizing the Causal Effect}
\label{sec:shapley_correction}
The direct Shapley formulation introduces a systematic scale difference between effects estimated for singleton and multi-concept groups. This section identifies the source of that distortion and develops a payout-coefficient correction intended to make the resulting causal effects comparable while retaining genuine interaction effects.

\subsection{Systematic Shrinkage of Joint Effects}
\label{sec:kl_ceiling}

Equations~\ref{eq:new_ce} and \ref{eq:new_ce_vector} technically complete the evaluation pipeline. Running an actual implementation, however, reveals a systematic pattern in the resulting scores. Whenever a concept $C_m$ is assigned to a correlation group $s$ with $|s|>1$, its resulting Shapley-CE tends to be lower than the causal effect that would have been measured had $C_m$ been evaluated on its own, i.e., on average $\mathrm{CE}(\mathbf{x},C_m,s) < \mathrm{CE}(\mathbf{x},C_m,\{C_m\})$, even for concepts that are otherwise independent in the sense of Section~\ref{sec:interaction_effects}.

The Shapley \emph{Efficiency} axiom (Section~\ref{sec:shapley_fundamentals}) makes the scale of the joint value relevant to every allocated effect: the sum of payouts to all concepts in $s$ equals $v_{\mathbf{x}}(s)$, the joint effect of intervening on the entire group at once. Whenever this joint value is smaller than the sum of the corresponding singleton values, the shortfall is distributed across the individual payouts $\mathrm{CE}(\mathbf{x},C_m,s)$. Empirically, this pattern occurs systematically on average in the considered setting, although its direction and magnitude may differ for individual groups, particularly when genuine redundancy or synergy is present.

Intuitively, probability space does not permit direct addition of separately normalized outputs. Suppose that intervening on Concept A alone shifts the model's predicted probability of outcome $b$ from $50\%$ to $80\%$, and that intervening on Concept B alone produces the same shift. Adding the two resulting probabilities would yield $160\%$, while adding the two isolated shifts to the baseline would yield $110\%$; neither is a valid probability. The actual joint intervention must instead produce a newly normalized distribution. This establishes non-additivity in probability space, but does not by itself determine whether the resulting KL-based joint effect is smaller or larger than the sum of singleton effects.

One structural source of this non-additivity is the Bayesian hierarchical modeling step presented by Matton et al. \cite{mattonWalkTalkMeasuring2024} in Appendix C.2 of their work. To reconstruct a probability distribution vector from sampled realizations of the answer variable $Y$, they employ a softmax function. Let $I$ denote the intervention condition under which responses are sampled, and let $\theta_{y\mid I}$ denote the latent logit assigned to answer option $y$ under that condition. Each scalar probability component is assembled as
\begin{equation}\label{eq:culprit_softmax}
	p_{y\mid I} = \frac{e^{\theta_{y\mid I}}}{\sum_{y' \in \mathcal{Y}} e^{\theta_{y'\mid I}}}, \qquad y \in \mathcal{Y},
\end{equation}
and the complete vector is $\mathbf{p}_{I}:=(p_{y\mid I})_{y\in\mathcal{Y}}$. Its key property is that the components are strictly confined to $[0,1]$ and sum to exactly $1$ across all possible predictions in $\mathcal{Y}$. This preserves the fundamental axiom that no set of mutually exclusive events can have a combined probability exceeding $100\%$.

Because of this strict normalization, the probability distribution resulting from a joint intervention cannot equal the sum of the distributions from individual interventions. If one were to naively add the isolated resulting distributions of an intervention on Concept A and an intervention on Concept B, the sum over the prediction space would equal $2$. Since any valid probability distribution must sum to exactly $1$, the joint intervention is inherently non-additive in probability space.

Beyond this extreme case of distributions summing to $2$, the nonlinear mapping from logits to probabilities means that additive changes in logit space need not remain additive in probability space. Example~\ref{ex:softmax_ceiling} of Appendix~\ref{app:examples} illustrates one sub-additive instance of this mechanism, in which two independent concepts that individually shift a prediction by $+20\%$ and $+30\%$ yield a joint shift of only $+44\%$, rather than the naive additive sum of $+50\%$. This example does not establish universal sub-additivity: depending on the intervention and on genuine concept interactions, a joint effect may also equal or exceed the sum of singleton effects.

Because the joint effect serves as the basis for calculating individual payouts, systematic non-additivity can create a reference-relative scale difference between singleton and multi-concept treatments. The observed average tendency $v_{\mathbf{x}}(s)<\sum_{C_m\in s}v_{\mathbf{x}}(\{C_m\})$ motivates a normalization, but is not imposed as a universal inequality. Rather than removing the softmax function or altering the Bayesian sampling process, the proposed normalization estimates the typical non-additivity of comparable treatments and adjusts against that reference level.

\subsection{Introducing Payout Coefficient}

For a treatment $T \subseteq s$ with $T \neq \emptyset$, the \emph{payout coefficient} is defined as
\begin{equation}
	\label{eq:payout_coefficient}
	\gamma(\mathbf{x},T) := \frac{v_{\mathbf{x}}(T)}{\displaystyle\sum_{C_m \in T} v_{\mathbf{x}}(\{C_m\})},
\end{equation}
the ratio of the group's joint effect to the sum of its members' singleton effects. If $\sum_{C_m\in T} v_{\mathbf{x}}(\{C_m\})=0$, every concept in $T$ has zero singleton effect; in this degenerate case $\gamma(\mathbf{x},T):=1$ is set by convention, so that no correction is applied.

By definition, $\gamma(\mathbf{x},T)=1$ indicates perfect additivity, $\gamma(\mathbf{x},T)<1$ the shrinkage documented above, and $\gamma(\mathbf{x},T)>1$ reinforcement. For a singleton treatment $T=\{C_m\}$, numerator and denominator coincide, so $\gamma(\mathbf{x},\{C_m\}) \equiv 1$: the correction defined below therefore leaves singleton effects unchanged, as required by \reqref{req:one_item_group}, introduced in Section~\ref{sec:requirements}.

\subsection{Constructing a Reference Set}
\label{sec:reference_set}

A treatment's own coefficient $\gamma(\mathbf{x},T)$ cannot be used to correct its own value function: doing so would force $v_{\mathbf{x}}(T)/\gamma(\mathbf{x},T) = \sum_{C_m\in T} v_{\mathbf{x}}(\{C_m\})$ identically for every treatment, discarding exactly the information about genuine concept interaction that \reqref{req:caputure_interaction}, introduced in Section~\ref{sec:requirements}, requires the metric to retain. Instead, $\gamma(\mathbf{x},T)$ is corrected using the \emph{typical} coefficient observed for treatments of the same size $|T|$ elsewhere in the dataset.

Let $\mathbf{S} := \mathcal{O}(\mathbf{x}, \mathbf{C}, n_{\max})$ denote the correlation groups for question $\mathbf{x}$. Three disjoint pools of reference treatments matching $|T|$ in size are considered, in descending priority:
\begin{align}
	\mathcal{R}_1(\mathbf{x},s,T) & := \big\{\, (\mathbf{x},T'') \;:\; T''\subseteq s,\ T''\neq T,\ |T''|=|T| \,\big\}, \label{eq:R_1}
	\\
	\mathcal{R}_2(\mathbf{x},s,T) & := \big\{\, (\mathbf{x},T'') \;:\; T''\subseteq s',\ s'\in \mathbf{S}\setminus\{s\},\ |T''|=|T| \,\big\}, \label{eq:R_2}
	\\
	\mathcal{R}_3(\mathbf{x},s,T) & := \big\{\, (\mathbf{x}'',T'') \;:\; \mathbf{x}''\in \mathbf{X}\setminus\{\mathbf{x}\},\ T''\subseteq s'',\ s''\in \mathcal{O}(\mathbf{x}'',\mathbf{C}(\mathbf{x}''),n_{\max}),\ |T''|=|T| \,\big\}. \label{eq:R_3}
\end{align}
$\mathcal{R}_1$ draws treatments of matching size from within the same correlation group $s$; $\mathcal{R}_2$ widens the search to other groups within the same question $\mathbf{x}$; $\mathcal{R}_3$ widens it further to the entire dataset $\mathbf{X}$. By construction, the three pools are pairwise disjoint. The reference set is obtained by drawing $n_{\mathrm{ref}} := 5$ elements greedily from $\mathcal{R}_1 \cup \mathcal{R}_2 \cup \mathcal{R}_3$, exhausting each pool before moving to the next:
\begin{equation}
	\label{eq:reference_set}
	\mathcal{R}(\mathbf{x},s,T) \subseteq \mathcal{R}_1\cup\mathcal{R}_2\cup\mathcal{R}_3, \qquad |\mathcal{R}(\mathbf{x},s,T)| = \min\big(n_{\mathrm{ref}},\ |\mathcal{R}_1\cup\mathcal{R}_2\cup\mathcal{R}_3|\big),
\end{equation}
with ties within a pool broken arbitrarily, e.g., uniformly at random. If $\mathcal{R}_1\cup\mathcal{R}_2\cup\mathcal{R}_3=\emptyset$, no reference treatment of matching size exists anywhere in the dataset, and no correction is applied (Section~\ref{sec:corrected_value_function}).

\subsection{Correcting the Causal Effect Value Function}
\label{sec:corrected_value_function}

The \emph{singleton-weighted reference coefficient} is formed by dividing the
aggregate joint effect of the reference treatments by their aggregate singleton
effect:
\begin{equation}
	\label{eq:gamma_bar}
	\bar\gamma(\mathbf{x},s,T) :=
	\begin{cases}
		\dfrac{\displaystyle\sum_{(\mathbf{x}'',T'')\in\mathcal{R}(\mathbf{x},s,T)}v_{\mathbf{x}''}(T'')}
		{\displaystyle\sum_{(\mathbf{x}'',T'')\in\mathcal{R}(\mathbf{x},s,T)}\sum_{C_i\in T''}v_{\mathbf{x}''}(\{C_i\})}
		  & \text{if the denominator is positive}, \\
		1 & \text{otherwise}.
	\end{cases}
\end{equation}
For non-degenerate reference treatments, this is equivalently a weighted average
of their payout coefficients, with weights given by their total singleton effects.
It therefore estimates the reference-relative shrinkage or amplification for
treatments of size $|T|$ without allowing references with negligible singleton
effects to receive the same influence as references carrying substantially more
effect mass. If no reference exists or the aggregate singleton denominator is
zero, the coefficient is set to $1$, so no normalization is applied. This
reference level also permits a distinction between two notions of independence.
\emph{Absolute additivity} means $\gamma(\mathbf{x},T)=1$, whereas
\emph{reference-relative independence} means
\begin{equation}
	\gamma(\mathbf{x},T)=\bar\gamma(\mathbf{x},s,T).
\end{equation}
The latter does not assert that the joint effect is absolutely additive; it states that the treatment exhibits the singleton-weighted level of non-additivity represented by equally sized reference treatments. Values below or above $\bar\gamma(\mathbf{x},s,T)$ indicate greater shrinkage or reinforcement, respectively, relative to that reference level.

To prevent division by a numerically degenerate reference coefficient, let $\varepsilon:=0.01$ denote a small fixed threshold. The corrected question-specific value function is
\begin{equation}
	\label{eq:v_adj}
	v_{\mathrm{adj},\mathbf{x}}(T) :=
	\begin{cases}
		\dfrac{v_{\mathbf{x}}(T)}{\bar\gamma(\mathbf{x},s,T)} & \text{if }T\neq\emptyset\text{ and }\bar\gamma(\mathbf{x},s,T)>\varepsilon,         \\
		v_{\mathbf{x}}(T)                                     & \text{if }T\neq\emptyset\text{ and }0\leq\bar\gamma(\mathbf{x},s,T)\leq\varepsilon, \\
		0                                                     & \text{if }T=\emptyset.
	\end{cases}
\end{equation}
Thus, no normalization is applied when the reference estimate is too close to zero to support stable division. The threshold is fixed before evaluation and reported with the experimental configuration.

Substituting $v_{\mathrm{adj},\mathbf{x}}$ for $v_{\mathbf{x}}$ in Equation~\ref{eq:new_ce} yields the corrected causal effect
\begin{equation}
	\label{eq:ce_adj}
	\mathrm{CE}_{\mathrm{adj}}(\mathbf{x},C_m,s) := \sum_{T\subseteq s\setminus\{C_m\}} \frac{|T|!\,(|s|-|T|-1)!}{|s|!} \big[v_{\mathrm{adj},\mathbf{x}}(T\cup\{C_m\}) - v_{\mathrm{adj},\mathbf{x}}(T)\big],
\end{equation}
which replaces $\mathrm{CE}(\mathbf{x},C_m,s)$ from this point onward in populating the causal-effect vector $\mathbf{CE}(\mathbf{x},\mathbf{C})$ of Equation~\ref{eq:new_ce_vector}.

In summary, the correction is designed to reduce the systematic, reference-relative scale distortion observed for multi-concept treatments while retaining deviations that may reflect interaction effects (Section~\ref{sec:interaction_effects}). Under the simplifying assumptions stated in Math Box~\ref{ma:dataset_average_preservation}, it is additionally intended to approximately preserve the original metric's dataset-average CE score.

\section{Remaining Evaluation Pipeline}
\label{sec:remaining_pipeline}
The previous sections introduced the modified formulation of the \emph{Causal Effect} (Equation~\ref{eq:ce_adj}), which is now computed using Shapley values. Furthermore, the underlying value function was extended to compensate for the systematic, reference-relative scale difference observed for joint concept effects and discussed in Section~\ref{sec:shapley_correction}.

After the CE vector has been computed, two major stages remain in the evaluation pipeline proposed by Matton et al.~\cite{mattonWalkTalkMeasuring2024}: estimating the \emph{Explanation Implied Effect} (EE) and computing the final faithfulness score from the relationship between both vectors. Fortunately, neither of these components requires any modification.

The computation of the EE vector remains unchanged because the output of the modified causal-effect estimation is still defined on a per-concept basis. Consequently, the adjusted CE vector is fully compatible with the inputs and outputs expected by the original EE estimation procedure. Although \emph{FELICS} explicitly models interactions between concepts, the objective of the EE vector remains identical: estimating which concepts the evaluated LLM considers influential enough to mention in its explanation.

For example, two concepts may each have only a minor individual influence on the prediction while jointly producing a substantial change due to a synergistic interaction. In this case, neither concept would be expected to appear frequently in explanations when considered individually, whereas both should be referenced once their combined effect becomes significant. Likewise, in a redundant setting, both concepts should still be mentioned whenever their joint contribution remains important for the prediction. Conversely, if one concept contributes only marginally to the coalition's overall effect, it should also appear less frequently in the explanation. Consequently, the alignment between causal effects and explanation-implied effects remains an appropriate indicator of explanation faithfulness.

Since the EE estimation remains unchanged, the final faithfulness estimation can likewise be adopted without modification. This applies both to the formal definition based on the \emph{Pearson Correlation Coefficient} (Equation~\ref{eq:faithfulness}) and to the practical estimation using \emph{Bayesian Hierarchical Modelling}, as proposed by Matton et al.~\cite{mattonWalkTalkMeasuring2024}. Furthermore, the computation of \emph{Dataset-Level Faithfulness} (Equation~\ref{eq:dataset_faithfulness}) is unaffected. By restricting all modifications to the estimation of causal effects, \emph{FELICS} builds upon the validated components of the original framework while minimizing the number of changes required throughout the remaining evaluation pipeline.

\section{Post-Modeling Assessment}
With the formal construction of \emph{FELICS} complete, this section consolidates the proposed changes and assesses their theoretical consequences. It compares the updated and original pipelines, revisits the requirements established earlier, and derives the effects expected from the new formulation before empirical evaluation.

\subsection{Reviewing Changes}

Having introduced the individual modifications of \emph{FELICS}, it is useful to summarize how the proposed framework differs from the original \emph{Causal Concept Faithfulness Metric} by Matton et al.~\cite{mattonWalkTalkMeasuring2024}. Rather than repeating the original evaluation pipeline presented in Section~\ref{sec:causal_concept_metric} and each modification introduced throughout this chapter, Figure~\ref{fig:comparison_matton_felics} provides a side-by-side comparison of both approaches.

The diagram is organized into three columns. The left column contains the original evaluation pipeline proposed by Matton et al. \cite{mattonWalkTalkMeasuring2024}, whereas the right column illustrates the corresponding components introduced by \emph{FELICS}. Between them, the shared components are shown to emphasize which parts of the framework remain unchanged. Whenever a processing step differs depending on the underlying dataset type (natural-language or tabular), the corresponding alternatives are separated by the \textbullet\ symbol. Finally, the arrows indicate how intermediate variables and results propagate through the evaluation pipeline and serve as inputs to subsequent processing stages.

\begin{figure}
	\centering
	\includestandalone[mode=buildnew, width=12.5cm]{graphs/comparison_matton_felics}
	\caption{Comparison of the evaluation pipeline between Matton et al.\ \cite{mattonWalkTalkMeasuring2024} (left) and the proposed \emph{FELICS} framework (right) with highlighted modifications.}
	\label{fig:comparison_matton_felics}
\end{figure}

\subsection{Reexamination of Requirements}

Section~\ref{sec:requirements} introduced six requirements that the proposed extension should satisfy. They cover preservation of the original metric's strengths (\reqref{req:preserving_strengths}), backward compatibility for singleton groups and independent concepts (\reqrefs{req:one_item_group}{req:independence_behavior}), the two central aspects of the \emph{Correlated Concepts Problem} (\reqrefs{req:no_surgical}{req:caputure_interaction}), and practical computational feasibility (\reqref{req:feasibility}).

At this stage, the conceptual design of \emph{FELICS} is complete, whereas its empirical evaluation using a reference implementation is presented in the following chapter. It is therefore appropriate to revisit these requirements and assess whether they are satisfied by the theoretical construction of the proposed framework.

\subsubsection{\ref{req:preserving_strengths} Preserving the strengths of the \emph{Causal Concept Faithfulness Metric}}

\reqref{req:preserving_strengths}, introduced in Section~\ref{sec:requirements}, requires the proposed extension to retain the principal strengths of the metric proposed by Matton et al.~\cite{mattonWalkTalkMeasuring2024}. These strengths stem from its operation on high-level semantic concepts rather than individual tokens and from its ability to estimate causal and explanation-implied effects for every concept separately. Since \emph{FELICS} retains this concept-based representation and continues to compute both the CE and EE vectors on a per-concept basis, these strengths are preserved without modification.

\subsubsection{\ref{req:one_item_group} Singleton-Group Compatibility}

\reqref{req:one_item_group}, introduced in Section~\ref{sec:requirements}, states that whenever every correlation group contains exactly one concept, the resulting causal-effect estimates should be identical to those obtained by Matton et al.~\cite{mattonWalkTalkMeasuring2024}. This follows directly from the properties of the normalization factor and the Shapley-value formulation, as demonstrated in Math Box~\ref{ma:one_item_group}.

\subsubsection{\ref{req:independence_behavior} Independence Compatibility}

\reqref{req:independence_behavior}, introduced in Section~\ref{sec:requirements}, considers the more general case in which concepts are grouped together but remain mutually independent. The nonlinear mapping from logits to probabilities discussed in Section~\ref{sec:kl_ceiling} means that even concepts without a substantive interaction need not satisfy $v_{\mathbf{x}}(s)=\sum_{C_m\in s}v_{\mathbf{x}}(\{C_m\})$ exactly.

The normalization proposed in Section~\ref{sec:shapley_correction} therefore uses \emph{reference-relative independence}: a correlation group is neutral relative to the dataset when its payout coefficient equals the singleton-weighted reference coefficient of equally sized treatments. This condition is distinct from absolute additivity, which would require a payout coefficient of exactly $1$. Under the additional reference-relative additivity assumptions stated in Math Box~\ref{ma:independence_behavior} of Appendix~\ref{app:mathematical_demonstrations}, no additional interaction effect is attributed to the group and the resulting causal-effect estimates coincide with those of the original metric.

\subsubsection{\ref{req:no_surgical} Avoiding the Assumption of Surgical Interventions}

\reqref{req:no_surgical}, introduced in Section~\ref{sec:requirements}, requires the proposed extension to avoid assuming that \emph{surgical interventions} are always possible. The original metric proposed by Matton et al.~\cite{mattonWalkTalkMeasuring2024} fundamentally relies on this assumption, although, as discussed in Section~\ref{sec:pratical_limit}, it does not hold in many practical scenarios where concepts are inherently correlated. Consequently, removing this assumption constituted one of the primary design goals of \emph{FELICS}.

Instead of intervening on individual concepts exclusively, \emph{FELICS} performs joint interventions whenever concepts belong to the same correlation group. This principle was introduced in Section~\ref{sec:joint_effect}. Assuming that the oracle correctly identifies concepts that cannot be meaningfully intervened upon in isolation, the framework no longer depends on the existence of surgical interventions. Therefore, \reqref{req:no_surgical} can be regarded as fulfilled.

\subsubsection{\ref{req:caputure_interaction} Correctly Capturing Interaction Effects}

\reqref{req:caputure_interaction}, introduced in Section~\ref{sec:requirements}, requires the proposed framework to distinguish between independent, redundant, and synergistic interactions and to attribute their causal effects accordingly, based on the interaction types established in Section~\ref{sec:interaction_effects}.

This property follows under the assumptions of Math Box~\ref{ma:independence_behavior}. If a correlation group exhibits reference-relative independence, its payout coefficient equals the singleton-weighted reference coefficient for treatments of identical size,
\begin{equation}
	\bar\gamma(\mathbf{x},s,s)=\gamma(\mathbf{x},s),
\end{equation}
causing the adjusted causal effects to reduce to those of the original metric. Conversely, if a correlation group exhibits stronger redundancy or synergy than comparable groups, its payout coefficient deviates from this reference value,
\begin{equation}
	\bar\gamma(\mathbf{x},s,s)\neq\gamma(\mathbf{x},s).
\end{equation}
As a consequence, the normalization either decreases or increases the adjusted coalition value before the Shapley-value attribution is performed. Redundant interactions therefore satisfy
\begin{equation}
	\mathrm{CE}_{\mathrm{adj}}(\mathbf{x},C_m,s)
	<
	\mathrm{CE}_{\mathrm{Matton}}(\mathbf{x},C_m),
\end{equation}
whereas synergistic interactions satisfy
\begin{equation}
	\mathrm{CE}_{\mathrm{adj}}(\mathbf{x},C_m,s)
	>
	\mathrm{CE}_{\mathrm{Matton}}(\mathbf{x},C_m).
\end{equation}
The proposed framework thus behaves consistently with the intended interpretation of interaction effects and can therefore be argued to satisfy \reqref{req:caputure_interaction}.

\subsubsection{\ref{req:feasibility} Maintaining Computational Feasibility}

\reqref{req:feasibility}, introduced in Section~\ref{sec:requirements}, requires the proposed extension to remain computationally feasible, since it would provide little practical benefit if its computational cost increased to an impractical level.

Introducing joint interventions naturally increases the number of counterfactual questions that must be generated. In the worst case, evaluating every possible concept combination would require an exponential number of LLM calls. To limit this growth, the oracle introduced in Section~\ref{sec:joint_counterfactual} groups together only those concepts for which joint interventions are considered necessary. Furthermore, the mandatory upper bound $n_{\max}$ ensures that the number of generated interventions never exceeds the limit established in Equation~\ref{eq:subset_counterfactual_bound}.

Whether \reqref{req:feasibility} is satisfied ultimately depends on what is considered computationally feasible in practice. Nevertheless, it can reasonably be argued that the proposed framework fulfills this requirement. Although \emph{FELICS} is necessarily more computationally expensive than the original metric of Matton et al.~\cite{mattonWalkTalkMeasuring2024}, this additional cost is restricted to questions containing correlated concepts. The combination of oracle-guided correlation detection and the hard upper bound on intervention size provides a principled trade-off between computational cost and modeling accuracy.

\subsection{Anticipated Effects}
\label{sec:influence}
Since the original framework by Matton et al.~\cite{mattonWalkTalkMeasuring2024} cannot reliably handle concepts acting in a redundant or synergistic manner, their true causal effect is not captured. The influence exerted by \emph{FELICS} on CE results relative to the \emph{Causal Concept Faithfulness} metric is determined by the intervention mode, as discussed in Section~\ref{sec:interaction_effects}.

\subsubsection{Effects on CE Scores}
The relationship between $\mathrm{CE}_{\text{adj}}$ and $\mathrm{CE}_{\mathrm{Matton}}$ is governed by the payout coefficient $\gamma$ relative to the singleton-weighted reference coefficient $\bar\gamma$. For a pair of concepts $C_i, C_j$ in a group $s = \{C_i, C_j\}$, where
\begin{equation}
	\gamma\bigl(\mathbf{x},\{C_i,C_j\}\bigr) = \frac{v_{\mathbf{x}}(\{C_i,C_j\})}{v_{\mathbf{x}}(\{C_i\}) + v_{\mathbf{x}}(\{C_j\})},
\end{equation}
three cases can be distinguished.

\noindent\textbf{Reference-relative redundancy ($\gamma(\mathbf{x},\{C_i,C_j\}) < \bar\gamma(\mathbf{x},s,s)$).} The joint effect is smaller relative to its singleton effects than for the reference treatments. When concepts are \emph{added} to a question, $\mathrm{CE}_{\text{adj}} < \mathrm{CE}_{\mathrm{Matton}}$ for both concepts, meaning that Matton et al.'s~\cite{mattonWalkTalkMeasuring2024} metric overestimates their effects. When concepts are \emph{removed}, $\mathrm{CE}_{\text{adj}} > \mathrm{CE}_{\mathrm{Matton}}$, meaning that the original metric underestimates them.

\noindent\textbf{Reference-relative synergy ($\gamma(\mathbf{x},\{C_i,C_j\}) > \bar\gamma(\mathbf{x},s,s)$).} The joint effect is larger relative to its singleton effects than for the reference treatments. When concepts are \emph{added}, $\mathrm{CE}_{\text{adj}} > \mathrm{CE}_{\mathrm{Matton}}$, so the original metric underestimates their effects. When concepts are \emph{removed}, $\mathrm{CE}_{\text{adj}} < \mathrm{CE}_{\mathrm{Matton}}$, and the original metric instead overestimates them.

\noindent\textbf{Reference-relative independence ($\gamma(\mathbf{x},\{C_i,C_j\}) = \bar\gamma(\mathbf{x},s,s)$).} The treatment exhibits the level of non-additivity typical of its references; this does not require absolute additivity. Under the compatibility assumptions stated in Math Box~\ref{ma:independence_behavior}, $\mathrm{CE}_{\text{adj}} = \mathrm{CE}_{\mathrm{Matton}}$ for both concepts, and FELICS matches Matton et al.'s~\cite{mattonWalkTalkMeasuring2024} original metric.

\subsubsection{Value-Replacement Interventions}
When concept values are \emph{modified} rather than added or removed, the relationship between $\mathrm{CE}_{\text{adj}}$ and $\mathrm{CE}_{\mathrm{Matton}}$ depends on the direction of the effect. If concepts $C_i$ and $C_j$ exhibit minor influence before intervention but major individual influence afterward while also interacting redundantly post-intervention, Matton’s metric overestimates their effect ($\mathrm{CE}_{\text{adj}} < \mathrm{CE}_{\mathrm{Matton}}$). Conversely, if their influence is major before intervention but minor after being changed on their own, again with redundant behavior after intervention, Matton’s metric underestimates ($\mathrm{CE}_{\text{adj}} > \mathrm{CE}_{\mathrm{Matton}}$). This sign inversion mirrors the behavior observed for addition vs. removal interventions.

However, value-based interventions introduce additional complexity: modifying a concept’s value may create contradictory information in the question (e.g., changing nationality from \texttt{Japanese} to \texttt{Indian} while retaining \texttt{Japanese} as the mother tongue). Such inconsistencies can obscure the measurement of causal effects, making it hard to predict the generalized impact of multi-concept modifications.

Despite these intervention-specific variations, the dataset-level average of all CE values is intended to remain approximately equal between both metrics under the simplifying assumptions stated in Math Box~\ref{ma:dataset_average_preservation} of Appendix~\ref{app:mathematical_demonstrations}.

\subsubsection{Effects on Faithfulness Scores}
As established in Section~\ref{sec:remaining_pipeline}, the definition of Explanation-Implied Effect (EE) remains entirely unchanged between FELICS and the \emph{Causal Concept Faithfulness} metric. Consequently, the EE vectors are identical for both metrics.

The faithfulness score, however, is still determined by the Pearson correlation coefficient between the CE and EE vectors. Since Pearson correlation depends on the magnitudes of the data points (unlike rank-based metrics such as Spearman), and the individual CE values may differ between FELICS and Matton’s metric, the resulting faithfulness scores are not guaranteed to be equal. In fact, whenever interactions are present, the alignment between CE and EE vectors generally changes, except in the rare case where widening gaps for some concepts are exactly offset by narrowing gaps for others. This behavior is illustrated in example~\ref{ex:payout_normalization_correlation} of Appendix~\ref{app:examples}.

Still, two special cases defined by \reqrefs{req:one_item_group}{req:independence_behavior} in Section~\ref{sec:requirements} yield identical faithfulness scores under their respective assumptions: when every correlation group contains exactly one concept and when every group satisfies the reference-relative additivity conditions stated in Math Box~\ref{ma:independence_behavior}.
